% !TeX root = ../main.tex
% Methodology section for a conference-length paper on Qwen3.5AE-4B ASR Stage-1.
% Scope: Stage-1 frozen-encoder ablation (5-encoder comparison: EnCodec, Mimi acoustic/semantic, DACVAE, Whisper-small).
% Stage-2 LoRA SFT (DACVAE-only v2 chain) is documented separately in docs/stage2/.
% One sentence per line; citation keys match methodology.md § 10.
% Target length: compact prose suitable for a single conference-paper section.
%
% Required packages (add to main preamble if missing):
%   \usepackage{amssymb}   % \mathbb{R}
%   \usepackage{booktabs}  % \toprule, \midrule, \bottomrule
%   \usepackage{amsmath}   % aligned math (align environment, \text{})
%   \usepackage{hyperref}  % optional, for clickable \cite / \ref
%   \usepackage[T1]{fontenc}  % optional, for | and < > rendering in verbatim/texttt

\section{Experiment Design}
% Encoder -> projector (Lamma) -> LM
% 각 encoder에서 feature 어디꺼 가져왔는지
% 학습 단계별 loss -> SFT
% 데이터셋, 각 encoder, evaluation benchmark, 분석 dataset

\label{sec:methodology}

\subsection{Architecture}
\label{sec:method:arch}

We adopt a three-stage encoder, projector, and LM architecture \cite{FLM25,SLAM24}, where a frozen audio encoder maps the waveform to a continuous latent sequence, a trainable projector aligns this latent with the LM embedding space, and a frozen LM autoregressively generates the transcript.

\paragraph{Audio encoders.}
We treat the audio encoder as a swappable component and study five pretrained candidates that span distinct architectural and training-objective design points, summarized in Table~\ref{tab:encoders}.
Four are neural audio codecs, namely \textbf{EnCodec}~\cite{Df22}, \textbf{Mimi}~\cite{Mi24} in its acoustic and semantic variants, and \textbf{DACVAE}~\cite{DAC23,AS24}, while the fifth is the supervised ASR encoder of \textbf{Whisper-small}~\cite{Wh23}.
For EnCodec, Mimi, and DACVAE we extract only the feed-forward encode path, omitting any decoder, residual vector quantizer, watermarker, or output projection; for Whisper we use the encoder stack only, discarding the text decoder.
In all cases the encoder is frozen throughout Stage-1 training.

\begin{table}[t]
\centering
\caption{Audio encoder candidates. ``Arch.'' abbreviates Conv (SEANet-style convolutional stack), Tr (Transformer), VAE (variational bottleneck). ``Obj.'' abbreviates Cdc (neural codec), Cdc\,+\,Sem (codec with a semantic-distillation Transformer head), Cdc\,+\,WM (codec with localized acoustic watermarking), ASR (supervised automatic speech recognition). Params are the parameter count of the active forward path; tokens/s $= \text{SR} / \text{hop}$. Pretrained checkpoint URIs are listed in Table~\ref{tab:encoder-checkpoints} of the appendix.}
\label{tab:encoders}
\small
\begin{tabular}{lllrrrrrr}
\toprule
Encoder            & Arch.     & Obj.         & SR (kHz) & Hop & Tokens/s & Dim. & Params & Ref. \\
\midrule
EnCodec            & Conv      & Cdc          & $24$ & $320$  & $75$ & $128$ & $7.43$M  & \cite{Df22} \\
Mimi (acoustic)    & Conv      & Cdc          & $24$ & $960$  & $25$ & $512$ & $12.63$M & \cite{Mi24} \\
DACVAE             & Conv+VAE  & Cdc + WM     & $48$ & $1920$ & $25$ & $128$ & $27.55$M & \cite{DAC23,AS24} \\
Mimi (semantic)    & Conv+Tr   & Cdc + Sem    & $24$ & $960$  & $25$ & $512$ & $37.82$M & \cite{Mi24} \\
Whisper-small      & Tr        & ASR          & $16$ & $320$  & $50$ & $768$ & $88.15$M & \cite{Wh23} \\
\bottomrule
\end{tabular}
\end{table}

\paragraph{Projector and LM.}
The LM is Qwen3.5AE-4B, an audio-enabled derivative of Qwen3.5-4B~\cite{QW24} with a hybrid linear and full attention backbone and reserved audio placeholder tokens.
The encoder output is mapped into the LM embedding space by a four-layer causal Llama decoder adapter~\cite{Ll23} whose total size stays within $\pm 2\%$ of $18.2$M parameters across all encoders in Table~\ref{tab:encoders}; see Appendix~\ref{app:projector} for details.
At forward time the projector output replaces the LM embeddings at positions marked by the placeholder token \texttt{<|audio\_pad|>}~\cite{QA24}, after which the LM generates the transcript autoregressively.

\subsection{Projector Alignment Objective}
\label{sec:method:objective}

Stage-1 training updates only the projector's $18.16$M parameters (roughly $0.45\%$ of the LLM) and freezes everything else.
Each sample is wrapped in the ChatML template~\cite{OAI23}: a system prompt, a user turn containing $t_{\text{audio}} = \lfloor n_{\text{samples}} / 1920 \rfloor$ copies of the audio placeholder token \texttt{<|audio\_pad|>} bracketed by \texttt{<|audio\_start|>} and \texttt{<|audio\_end|>}~\cite{QA24}, a fixed instruction ``Transcribe the audio to text.'', and an assistant turn containing the ground-truth transcript followed by the end-of-sequence token.
At forward time the projector output replaces the \texttt{<|audio\_pad|>} embeddings in the LLM input, producing an input sequence whose audio-conditioned region has the same length as the projector output.
The training objective is standard causal cross-entropy, but restricted to the transcript and EOS tokens: all system, user, audio-placeholder, and instruction tokens are marked with the ignore-index of $-100$~\cite{HF23}.
Further training details, including sequence packing, optimizer settings, learning-rate schedule, and system-level kernels, are deferred to Appendix~\ref{app:training}.

\subsection{Data}
\label{sec:method:data}

We train on an English ASR mixture of LibriTTS-R~\cite{Ko23}, the English subset of Multilingual LibriSpeech~\cite{Pr20}, and the English portion of VoxPopuli~\cite{Wa21}, totalling $46{,}074$ hours (Table~\ref{tab:datasets}).
Audio is streamed on demand from an internal object storage; files stored at $16$\,kHz are resampled to $48$\,kHz online to match the DACVAE sampling rate.
Starting from a pre-shuffled 128-shard JSONL manifest that also contained CommonVoice and GigaSpeech, we retain only the three corpora in Table~\ref{tab:datasets} by substring-matching the audio path, yielding $11{,}345{,}224$ entries ($54.9\%$ of the original) across $128$ roughly equal shards.
The streaming pipeline is built on Hugging Face \texttt{IterableDataset}~\cite{HF23}, with shards distributed across data-parallel ranks via \texttt{split\_dataset\_by\_node}, a per-rank packing buffer of $128$ utterances, a $45$-second maximum audio length, and a $12$-worker download pool.

\begin{table}[t]
\centering
\caption{Stage-1 English ASR training data.}
\label{tab:datasets}
\small
\begin{tabular}{lrl}
\toprule
Corpus & Duration (h) & Ref. \\
\midrule
LibriTTS-R           & $552.42$      & \cite{Ko23} \\
MLS (en)             & $45{,}000.00$ & \cite{Pr20} \\
VoxPopuli (en)       & $522.00$      & \cite{Wa21} \\
\midrule
Total                & $46{,}074$    &             \\
\bottomrule
\end{tabular}
\end{table}

\subsection{Evaluation Benchmark}
\label{sec:method:eval}

% TODO: finalize benchmark selection.
% Candidate benchmarks to consider:
%   - LibriSpeech test-clean / test-other (word error rate)
%   - TED-LIUM 3 (out-of-domain read speech)
%   - VoxPopuli test (European Parliament, conversational)
%   - CommonVoice test (accented / noisy, cross-domain)
%   - GigaSpeech test (in-the-wild podcast / YouTube)
%   - FLEURS (multilingual — English subset for zero-shot cross-lingual check)
% TODO: decide primary metric. Default = WER after Whisper-style text normalization.
% TODO: decide decoding strategy. Default = greedy with temperature 0; also report beam search if competitive baselines require it.
% TODO: decide whether to report both in-domain and out-of-domain results separately.

We evaluate all encoder variants on a common set of English ASR benchmarks that cover in-domain read speech, out-of-domain read speech, and spontaneous / conversational speech.
\emph{[Benchmark list TBD.]}
The primary metric is word error rate (WER) after Whisper-style text normalization~\cite{Wh23}, and we additionally report character error rate (CER) for fine-grained comparison.
Decoding is performed greedily with temperature $0$ and no additional language model.
All encoder variants share identical evaluation code, tokenization, and normalization to isolate the effect of the audio representation.

\subsection{Analysis Settings and Methods}
\label{sec:method:analysis}

% TODO: finalize analysis design.
% Candidate axes / probes:
%   - Paralinguistic probing: prosody, emotion, speaker attributes (see FLM25, LSTN26).
%   - Robustness: noise injection, speed perturbation, volume scaling.
%   - Representation geometry: CKA / linear probing across encoder layers.
%   - Token-rate ablation: downsample encoder output to match a common FPS.
%   - Dead-weight cost: forward latency / GPU memory per encoder.
%   - Sample-efficiency: WER vs. training-step curves.
% TODO: choose datasets for each analysis (may reuse subsets of the training / eval data).
% TODO: decide whether analysis is per-encoder ablation, or joint comparison.

Beyond benchmark WER, we analyse each encoder along several axes intended to expose how its pretraining objective and architecture shape the downstream ASR behaviour.
\emph{[Analysis axes and datasets TBD.]}
Candidate axes under consideration include (i) paralinguistic sensitivity~\cite{FLM25,LSTN26}, (ii) robustness to additive noise and speed perturbation, (iii) representation geometry across encoder layers measured by centered kernel alignment and linear probing, (iv) the effect of matching all encoders to a common token rate by downsampling, and (v) wall-clock and memory cost of the encoder forward pass.
All analyses share the Stage-1 projector weights and differ only in the evaluation-time inputs or probes, so they do not require retraining.

